\begin{abstract}

\textbf{Background:}
Variant effect prediction (VEP) is a bottleneck in clinical variant
interpretation. Existing tools are restricted to specific variant types---missense
(AlphaMissense, REVEL), splice-altering (SpliceAI)---or incorporate clinical
databases in training (CADD), creating potential circularity. No single tool
provides unified, training-data-free predictions across coding, non-coding, and
indel variants. Spaceflight radiation generates characteristic transversion
mutations and complex DNA damage in genes critical for radiation response,
telomere maintenance, and DNA repair, yet comprehensive variant effect maps
for these genes are lacking.

\textbf{Results:}
We applied Evo2, a 7-billion-parameter genomic foundation model trained on 9.3
trillion nucleotides, to score 215,001 variants (168,409 SNVs and 46,592 indels)
across ten spaceflight radiation-response genes in a zero-shot manner requiring
no training data. Deep mutational scanning calibration against four independent
datasets yielded Spearman $|\rho| = 0.21$--$0.52$. ClinVar benchmarking
demonstrated AUROCs of 0.91--1.00 across six genes, with Evo2 achieving the
highest AUROC for CHEK2 (0.9996) and TP53 (0.9887). On matched variant sets,
Evo2 was statistically indistinguishable from CADD (DeLong $p > 0.08$, all
genes). Evo2 scores were orthogonal to supervised tools (Spearman $\rho =
0.41$--$0.83$ vs.\ CADD), and an Evo2+CADD ensemble improved classification
by up to 0.38\%. For non-coding variants, Evo2 was the only tool with a
significant positive correlation with TERT promoter MPRA functional data
($\rho = +0.267$, $p = 3.8 \times 10^{-14}$; CADD, SpliceAI, ncER: not
significant). Frameshift indel AUROCs exceeded 0.94 across six genes (mean
0.977). Transversion-control analysis revealed that the elevated scores for
radiation-signature mutations are attributable to general transversion bias
rather than radiation-specific chemistry.

\textbf{Conclusions:}
Evo2 provides a comprehensive zero-shot variant effect atlas that complements
supervised tools and fills critical gaps for non-coding and indel variants.
The atlas of 215,001 scored variants across ten genes is provided as a
public resource for clinical variant interpretation and astronaut genomic
risk assessment.

\end{abstract}

\keywords{variant effect prediction, genomic foundation model, Evo2, deep
mutational scanning, ClinVar, spaceflight radiation, non-coding variants,
zero-shot prediction, TERT, DNA damage response}
